\section{从响应核到可辨识参数}
\label{chap:feqo-end-to-end}

前面的谱与收集公式确定了理想信号。实验反演还需要明确未知量、仪器响应和计数模型。下面把这些层次逐一连接，并用可计算的核空间检验哪些参数确实能从数据确定。

\begin{definition}[前向模型与参数层次]
用 $x$ 表示待估对象，用 $\mathcal F$ 表示给定电子源、几何、频窗和探测器后的前向映射：
\begin{equation}
d=\mathcal F(x)+\epsilon.
\label{eq:feqo-forward-operator}
\end{equation}
可以选择响应核、场分布或有限维材料参数作为 $x$，但这几种选择的线性性不同。固定轨迹时，EELS 对响应核 $\operatorname{Im}\mathbf G^R$ 是线性的；响应核对介电函数和几何通常非线性。CL 需要复 Green 张量到收集通道的振幅，PINEM 需要外加相干场。
\end{definition}

\begin{proposition}[三种读出的前向泛函]
对于本章的互易介质，沿 $z$ 的匀速电子在适用的弱损失条件下满足
\begin{equation}
\Gamma_{\rm EELS}(\omega;\vb{R})
=\frac{\mu_0q_{\mathrm e}^2}{\pi\hbar}
\int\dd z\dd z'\,\ee^{\ii\omega(z-z')/v}
\operatorname{Im}G^R_{{\rm sc},zz}
((\vb{R},z),(\vb{R},z');\omega).
\label{eq:feqo-eels-forward-green}
\end{equation}
互易情形下对调 $z,z'$ 给出相同实数。固定电子源以后，核的映射为
\begin{equation}
\mathcal F_{\rm EELS}:\operatorname{Im}\mathbf G^R_{\rm sc}\longmapsto\Gamma_{\rm EELS}.
\label{eq:feqo-eels-forward-map}
\end{equation}
对于具有通常真空远场展开的有限结构，定义
\[
\mathbf G_\infty(\hat{\vb{n}},\vb{r}';\omega)
=\lim_{r\to\infty}r\ee^{-\ii kr}\mathbf G^R(r\hat{\vb{n}},\vb{r}';\omega),
\quad
\vb{F}=\ii\mu_0\omega\int\mathbf G_\infty\widetilde{\vb{J}}_e\dd^3r'.
\]
取 $\vb{E}(t)=(2\pi)^{-1}\int\widetilde{\vb{E}}(\omega)\ee^{-\ii\omega t}\dd\omega$。若 $\mathsf Q(\hat{\vb{n}},\omega)$ 是偏振空间中满足 $0\le\mathsf Q\le I$ 的接收效率矩阵，则
\begin{equation}
\Gamma_{\rm CL}^{\rm det}
=\frac{1}{\pi Z_0\hbar\omega}
\int\dd\Omega\,\vb{F}^\dagger\mathsf Q\vb{F},
\qquad Z_0=\sqrt{\mu_0/\varepsilon_0}.
\label{eq:feqo-cl-forward-collection}
\end{equation}
由此
\begin{equation}
\mathcal F_{\rm CL}:\mathbf G_\infty\longmapsto\Gamma_{\rm CL}^{\rm det}.
\label{eq:feqo-cl-forward-map}
\end{equation}
只给出 $\operatorname{Im}G_{{\rm sc},zz}$ 不能确定这一通道振幅。对外加场，沿用本书正频复幅定义，有
\begin{align}
\beta(\vb{R})&=\frac{q_{\mathrm e}}{\hbar\omega}
\int E^{(+)}_z(\vb{R},z)\ee^{-\ii\omega z/v}\dd z,
\label{eq:feqo-pinem-forward-beta}\\
\mathcal F_{\rm PINEM}:E^{(+)}_z&\longmapsto
\{J_\ell^2(2|\beta|)\}_{\ell\in\ZZ}.
\label{eq:feqo-pinem-forward-map}
\end{align}
最后一式使用窄带、无反冲能量格和经典相干驱动条件。
\end{proposition}
\begin{proof}
第一式由电子电流代入\eqnrefs{eq:eels_green_master_dp47}，使用 $\dd z=v\dd t$ 得到。远场功率为 $r^2|\vb{E}|^2/Z_0$，正频 Parseval 公式给系数 $1/\pi$；除以单光子能量 $\hbar\omega$ 并投影接收偏振，得到 CL 式。PINEM 的线积分与边带概率则分别由\eqnrefs{eq:pinem_beta_green_unified_dp47}及经典驱动解得到。
\end{proof}

\begin{workedexample}[有限窗口核的伴随]
固定一个频率，将轨迹截在 $[-L/2,L/2]$，未知核取 $H\in L^2([-L/2,L/2]^2)$。记 $C=\mu_0q_{\mathrm e}^2/(\pi\hbar)$，则
\[
TH=C\int\dd z\dd z'\,\ee^{\ii\omega(z-z')/v}H(z,z'),
\qquad |TH|\le |C|L\|H\|_2.
\]
求该前向算子的伴随，并说明有限轨迹模型为何给出有界映射。

\textbf{解。} 取内积 $\langle f,g\rangle=\int f^*g$，有
\[
(T^*y)(z,z')=C^*\ee^{-\ii\omega(z-z')/v}y.
\]
这才是一个有界算子的伴随。在整个六维体积 $L^2$ 空间中，“把核限制到理想零宽轨迹”并不自动有界；实际束宽、空间正则性或有限维离散化必须进入模型。多位置观测对应多个这样的轨迹核，其伴随为各核的加权和。
\end{workedexample}

\subsection{参考相位、仪器响应与计数}

\begin{workedexample}[高斯近场的三相位复振幅读取]
设某条轨迹上的正频场为
\[
E_z^{(+)}(z)=E_*\ee^{\ii\phi}\ee^{-z^2/(2d^2)}\ee^{\ii k_z z},
\qquad\Delta k=k_z-\omega/v.
\]
求该轨迹上的复耦合，并用已标定参考的三个相位恢复它的两个正交分量；比较弱耦合计数与精确边带计数。

\textbf{解。} 直接完成高斯积分，得到
\begin{equation}
\beta_s=\frac{q_{\mathrm e}E_*\sqrt{2\pi}d}{\hbar\omega}
\ee^{-d^2\Delta k^2/2}\ee^{\ii\phi}
\equiv A\ee^{\ii\varphi}.
\label{eq:feqo-imaging-beta}
\end{equation}
$A\ge0$；电子电荷 $q_{\mathrm e}=-e$ 的负号并入 $\varphi$。沿同一相干电子轨迹叠加已标定参考耦合 $B\ee^{\ii\theta}$，则 $\beta_\theta=\beta_s+B\ee^{\ii\theta}$。记 $r_\theta=|\beta_\theta|$，弱耦合时
\begin{equation}
P_{+1}=P_{-1}=r_\theta^2+O(r_\theta^4),
\qquad P_0=1-2r_\theta^2+O(r_\theta^4).
\label{eq:feqo-imaging-weak-forward}
\end{equation}
写 $\beta_s=x+\ii y$，理想弱耦合强度 $Y_\theta=r_\theta^2$ 满足
\[
Y_\theta=x^2+y^2+B^2+2B(x\cos\theta+y\sin\theta).
\]
对 $\theta=0,\pi/2,\pi$，令 $C_0=(Y_0+Y_\pi)/2$，便有
\[
x=\frac{Y_0-Y_\pi}{4B},\qquad
y=\frac{Y_{\pi/2}-C_0}{2B},\qquad B>0.
\]
取 $A=B=0.05$、$\varphi=\pi/3$，三次测量的 $Y_\theta$ 为
\[
0.0075,\qquad0.005+0.0025\sqrt3\simeq0.00933013,
\qquad0.0025.
\]
每个相位入射 $10^5$ 个电子时，单侧一级边带的弱耦合期望计数约为 $750,933,250$。精确拟合应使用 $J_1^2(2r_\theta)$；相应计数约为 $744,924,249$，体现截断误差已经可见。
\end{workedexample}

\begin{definition}[归一化探测响应与计数模型]
若一级边带间距为 $\hbar\omega$，第 $\ell$ 条真实谱线经高斯分辨函数落入像素窗 $\mathcal B_j$ 的概率为
\[
R_{j\ell}=\int_{\mathcal B_j}\frac{\dd E}{\sqrt{2\pi}\sigma_E}
\exp\!\left[-\frac{(E-E_0-\ell\hbar\omega)^2}{2\sigma_E^2}\right].
\]
加入未检测结果后，每列满足 $\sum_jR_{j\ell}=1$。相位 $\theta_s$ 下的像素概率与 Poisson 入射模型为
\begin{equation}
p_{js}=\sum_\ell R_{j\ell}J_\ell^2(2|\beta_{\theta_s}|),
\qquad\mu_{js}=\nu_s p_{js}+b_{js},\qquad
n_{js}\sim\operatorname{Poisson}(\mu_{js}).
\label{eq:feqo-imaging-detector}
\end{equation}
独立 Poisson 计数适用于独立 Poisson 入射经分流的模型。若入射总数固定，各出口计数服从多项分布，有负协方差；不能把固定总数数据直接当作独立 Poisson 样本。
\end{definition}

\begin{proposition}[局部可辨识性的 Fisher 判据]
设有限维实参数 $\vb{\vartheta}$ 的所有 $\mu_j>0$，导数存在，且采用独立 Poisson 模型。令 $J_{ja}=\partial_a\mu_j$，则
\begin{equation}
I=J^{\mathsf T}\operatorname{diag}(\mu_j^{-1})J,
\qquad\operatorname{Cov}(\widehat{\vb{\vartheta}})\succeq I^{-1}.
\label{eq:feqo-cramer-rao}
\end{equation}
后一不等式要求 $I$ 可逆、估计无偏且满足通常微分积分互换条件。$I$ 正定当且仅当加权雅可比矩阵满列秩。
\end{proposition}
\begin{proof}
分数函数为 $\partial_a\log\mathcal L=\sum_j(n_j/\mu_j-1)\partial_a\mu_j$。利用 $\operatorname{Var}n_j=\mu_j$ 和独立性得到 Gram 矩阵 $I$。对任意向量 $h$，
\[
h^{\mathsf T}Ih=\sum_j\frac{(J_jh)^2}{\mu_j}.
\]
因此存在不可辨识的一阶方向当且仅当 $Jh=0$。Cram\'er--Rao 界由估计误差与分数函数的协方差恒等式及矩阵 Cauchy--Schwarz 不等式得到。
\end{proof}

\begin{figure}[htbp]
\centering
\includegraphics[width=.82\linewidth]{electron_field_calibration.pdf}
\caption{相位失配压低有效耦合，完整 Bessel 分布决定可见边带。高阶边带显著时，弱耦合近似不能用于定量反演；场幅估计还需校准轨迹积分与仪器响应。}
\label{fig:electron-field-calibration-modern}
\end{figure}

\subsection{正则化、噪声放大与空间分辨}

\begin{theorem}[有限维 Tikhonov 解与唯一性]
对白化并统一量纲后的数据 $d\in\CC^m$、未知系数 $x\in\CC^n$，给定 $A,L$ 和 $\eta>0$，定义
\begin{equation}
\widehat x_\eta=\mathop{\rm argmin}_x
\{\|Ax-d\|_2^2+\eta\|Lx\|_2^2\}.
\label{eq:feqo-tikhonov-regularization}
\end{equation}
解唯一当且仅当 $\ker A\cap\ker L=\{0\}$；此时
\begin{equation}
(A^\dagger A+\eta L^\dagger L)\widehat x_\eta=A^\dagger d.
\label{eq:feqo-tikhonov-normal}
\end{equation}
\end{theorem}
\begin{proof}
对实部和虚部求变分得到正规方程。二次型满足
\[
x^\dagger(A^\dagger A+\eta L^\dagger L)x
=\|Ax\|_2^2+\eta\|Lx\|_2^2,
\]
故正定性等价于公共核仅含零向量；若存在公共零方向，目标函数沿该方向不变。
\end{proof}

\begin{proposition}[奇异值滤波与误差界]
取 $L=I$，$A$ 的秩为 $r$，非零奇异值与奇异向量为 $\sigma_j,u_j,v_j$，$1\le j\le r$。真实数据 $d=Ax_*+\epsilon$、$\|\epsilon\|\le\delta$，$P_{\ker A}$ 为未知空间中核的正交投影。则
\begin{equation}
\widehat x_\eta=\sum_{j=1}^r\frac{\sigma_j}{\sigma_j^2+\eta}
\langle u_j,d\rangle v_j,
\label{eq:feqo-imaging-field-inversion}
\end{equation}
并且
\begin{align}
\|\widehat x_\eta-x_*\|
&\le\frac{\delta}{2\sqrt\eta}+\|b_\eta\|,\label{eq:feqo-tikhonov-error-bound}\\
b_\eta&=P_{\ker A}x_*+
\sum_{j=1}^r\frac{\eta}{\sigma_j^2+\eta}\langle v_j,x_*\rangle v_j.
\label{eq:feqo-tikhonov-bias}
\end{align}
第一项偏差 $P_{\ker A}x_*$ 完全不可从数据恢复。若额外有源条件 $x_*=A^\dagger z$，该项为零，总偏差不超过 $\sqrt\eta\|z\|/2$。
\end{proposition}
\begin{proof}
逐个奇异向量求正规方程得到滤波式。标量函数 $\sigma/(\sigma^2+\eta)$ 的最大值为 $1/(2\sqrt\eta)$，控制噪声放大；同一函数乘 $\eta$ 控制源条件下的偏差。没有源条件时，偏差必须保留其显式谱分解。
\end{proof}

\begin{workedexample}[高斯模糊与边带信息的秩]
设仪器空间响应为宽度 $w$ 的归一化高斯核，未知场离散为 $n$ 个复系数。求噪声为 $\delta_k$ 时可分辨的空间频率条件，以及固定轨迹、固定频率 PINEM 计数能够约束的实参数数目。

\textbf{解。} 一维高斯空间模糊的 Fourier 奇异值为 $\sigma(k)=\ee^{-w^2k^2/2}$。某一空间 Fourier 分量只有在 $\sigma(k)|x_k|\gtrsim\delta_k$ 时可分辨，因此信噪比大于一时满足
\[
|k|\lesssim\frac1w\sqrt{2\log(|x_k|/\delta_k)}.
\]
这个判据同时依赖仪器宽度和该分量的信噪比。

\begin{samepage}
在固定轨迹、固定频率下，任意多个 PINEM 边带的概率都只通过一个复数 $\beta[E]$ 依赖场。若离散场有 $n$ 个复系数，计数对其 $2n$ 个实坐标的雅可比矩阵必须因式分解为
\[
J_{\rm counts,field}=J_{\rm counts,\beta}J_{\beta,\rm field},
\qquad\operatorname{rank}J_{\rm counts,field}\le2.
\]
\begingroup\interlinepenalty=10000
无参考时还只依赖 $|\beta|$，局部秩至多一。高阶边带改善同一耦合的统计估计；增加独立空间核需改变轨迹、频率、入射场或观测通道。\par
\endgroup
\end{samepage}
\end{workedexample}
